\documentclass[10pt]{beamer}

\usepackage{amsmath,amsthm,amssymb,bm}
\usepackage{xcolor}
\usepackage{cancel}
\usepackage{graphicx}
\usepackage{changepage}
\usepackage{circuitikz}
\usepackage{pgfplots}
\usepackage{physics}
\usepackage{hyperref}
\usepackage{siunitx}
\usepackage{minted}
\usepackage{fontspec}
\usepackage{relsize}
\usepackage{subfig}
\usepackage{mhchem}
\usepackage{todonotes}
\usepackage{biblatex}
\usepackage{multicol, multirow, booktabs}
\usepackage[breakable]{tcolorbox}
% \usepackage[inline]{enumitem}

\pgfplotsset{compat=newest}
\usetikzlibrary{lindenmayersystems}
\usetikzlibrary{arrows}
\usetikzlibrary{calc}
\usetikzlibrary{positioning, fit}
\usetikzlibrary{3d, perspective}

\definecolor{incolor}{HTML}{303F9F}
\definecolor{outcolor}{HTML}{D84315}
\definecolor{cellborder}{HTML}{CFCFCF}
\definecolor{cellbackground}{HTML}{F7F7F7}
\newcommand{\ui}{\hat{i}}
\newcommand{\uj}{\hat{j}}
\newcommand{\uk}{\hat{k}}
\newcommand{\ux}{\hat{x}}
\newcommand{\uy}{\hat{y}}
\newcommand{\uz}{\hat{z}}
\newcommand{\uv}[1]{\hat{\bm{{#1}}}}
\newcommand{\pr}[1]{{#1^\prime}}
\newcommand{\justif}[2]{&{#1}&\text{#2}}
\newcommand{\dul}[1]{\underline{\underline{#1}}}
\newcommand{\pdiff}[2][]{{\frac{\partial^{#1}}{\partial {{#2}^{#1}}}}}
% \newcommand{\dif}[2][]{{\frac{d^{#1}}{d{{#2}^{#1}}}}}
\pgfdeclarelayer{background}  
\pgfsetlayers{background,main}
\AtBeginDocument{\RenewCommandCopy\qty\SI}

\makeatletter
\newcommand{\boxspacing}{\kern\kvtcb@left@rule\kern\kvtcb@boxsep}
\makeatother
\newcommand{\prompt}[4]{
    \ttfamily\llap{{\color{#2}[#3]:\hspace{3pt}#4}}\vspace{-\baselineskip}
}

\newcommand{\thevenin}[2]{
  \begin{center}
    \begin{circuitikz} \draw
      (0,0) -- (2,0) to[battery1, l_=$V_{Th}\eq#1$] (2,2) 
      to[resistor, l_=$R_{Th}\eq#2$] (0,2)
      ;
      \draw [o-] (-.07,2.079);
      \draw [o-] (-.07,0.079);
    \end{circuitikz}
  \end{center}
}

\newcommand{\norton}[2]{
  \begin{center}
    \begin{circuitikz} \draw
      (0,0) -- (3,0) to[american current source, l_=$I_{N}\eq#1$] (3,2) -- (0,2) (2,0)
      to[resistor, l=$R_{N}\eq#2$] (2,2)
      ;
      \draw [o-] (-.07,2.079);
      \draw [o-] (-.07,0.079);
    \end{circuitikz}
  \end{center}
}

\newcommand{\highlight}[1]{\colorbox{yellow}{#1}}

\newcommand{\ti}[1]{\widetilde{#1}}

\newfontface{\Kaufmann}{Kaufmann}
\DeclareTextFontCommand{\kf}{\Kaufmann}
\newcommand{\scriptr}{\fontsize{12pt}{12pt}\kf{r}}

\newfontface{\KaufmannB}{Kaufmann Bold}
\DeclareTextFontCommand{\kfb}{\KaufmannB}
\newcommand{\bscriptr}{\fontsize{12pt}{12pt}\kfb{r}}

\newcommand{\bv}[1]{\bm{\mathrm{#1}}}

\title{Physics 4605H: Project Presentation}
\author{Jeremy Favro (0805980) \\ Trent University, Peterborough, ON, Canada}

\begin{document}

\begin{frame}
  \titlepage
\end{frame}

\begin{frame}
  \frametitle{What we're doing}
  \begin{itemize}
    \item We have an $N\times N$ matrix $\dul{M}$ and a vector $\bv{b}$ and want to solve the equation $\dul{M}\bv{x}=\bv{b}$.
    \item The time complexity scales at least as order $N$ ($\order{N}$), this is how long it takes to write the solutions!
    \item Using a quantum algorithm we can determine the error in the fit without knowing parameters
  \end{itemize}
\end{frame}

\begin{frame}
  \frametitle{Why we're doing it: least squares}
  \begin{itemize}
    \item Fitting is the process of finding a nice $f(x)=\sum_j \phi_j(x)\lambda_j$ which approximates a dataset $\left\{(x_i,y_i)\right\}$
    \item Least squares is an algorithm to do this by minimizing $\sum_i\abs{f(x_i)-y_i}^2$
    \item This can be reduced to minimizing $\abs{\dul{F}\bv{\lambda}-\bv{y}}^2$ for
          \begin{itemize}
            \item $\dul{F}_{ij}=\phi_j(x_i)$
            \item $\bv{\lambda}$ is a coefficient vector for $\left\{\lambda_j\right\}$
            \item $\bv{y}$ is the goal data
          \end{itemize}
    \item The matrix must be \emph{sparse}
  \end{itemize}
\end{frame}

\begin{frame}
  \frametitle{How we're doing it: matrix math}
  \begin{itemize}
    \item We want to solve $\dul{F}\bv{\lambda}=\bv{y}$
    \item We apply the \emph{Moore-Penrose pseudoinverse} (MPPI), $$\dul{F}^+=(\dul{F}^\dagger\dul{F})^{-1}F^\dagger,$$ to both sides, get $\lambda=\dul{F}^+\bv{y}$
    \item With this $\dul{F}$ need not be square/invertible
    \item So to ``go quantum'' we need to find $\dul{F}^+$ and apply it
          % \item \todo{Talk about MPPI in speaker notes. What is it? Uses?}
  \end{itemize}
\end{frame}

\begin{frame}
  \frametitle{How we're doing it: going quantum}
  \begin{itemize}
    \item We need an input state corresponding to $\bv{y}$ and the ability to copy this state (an \emph{oracle})
    \item Also need oracles which give us states corresponding to entries in $\dul{F}$ and $\dul{F}^\dagger$
    \item Then we need to construct the pseudoinverse and apply it to the input $\bv{y}$
    \item This is hard, mathematically and to implement on a QC at all
    \item It gives a time complexity of $\order{\log(N)(\frac{s^3\kappa^6}{\epsilon})}$ where $\epsilon$ is the fit error tolerance
  \end{itemize}
\end{frame}

\begin{frame}
  \frametitle{What we're doing: error estimation}
  \begin{itemize}
    \item Because this is hard, it would be nice to know if it is viable
    \item Classical computers need to fit to determine fit \emph{error}. Takes $\order{N}$ as before
    \item A quantum algorithm can determine fit error/accuracy without fitting
  \end{itemize}
\end{frame}

\begin{frame}
  \frametitle{How we're doing it: swap test}
  \begin{itemize}
    \item We need the same oracles (magic boxes) as before, and a tolerance within which to determine error
    \item We perform the \emph{swap test} by
          \begin{itemize}
            \item Apply a controlled swap to states representing $\ket{\bv{y}}$ and $\dul{F}\ket{\bv{\lambda}}$ controlled by a qubit in superposition
            \item Apply a Hadamard gate to the superposition control qubit
            \item The states are different if measuring the control returns a 1. This happens with a calculable probability
            \item From the probability the error can be found
          \end{itemize}
  \end{itemize}
\end{frame}


\begin{frame}
  \frametitle{What we've done}
  \begin{itemize}
    \item Wanted to solve linear systems so we can fit a function to data
    \item Sketched a way to do so with a quantum computer \& the time complexity of doing so
    \item Implementing this is hard, so sketched a way to determine the error in a fit without fitting
  \end{itemize}
\end{frame}
\end{document}